\documentclass[11pt]{article}
\usepackage[margin=1.2in]{geometry}
\usepackage{amsmath,amssymb,amsthm}
\usepackage{hyperref}
\hypersetup{colorlinks=true, linkcolor=blue, urlcolor=blue}
\usepackage{parskip}
\emergencystretch=3em

\newcommand{\R}{\mathbb{R}}
\newcommand{\code}[1]{\texttt{#1}}

\title{Tide retrospective: the forward expansion domain\\
\large the Peano mixin and the exact exponent split}
\author{Claude (Anthropic), with GPT-5.6 Sol consults}
\date{2026-08-10}

\begin{document}
\maketitle

\begin{abstract}
The order-$N$ forward moment expansion needs exactly one hypothesis
beyond the merged higher-order Laplace package: the Peano form of the
Taylor remainder. The quantitative $O(\|y\|^{N+2})$ bound already in
\code{HigherLaplaceDomain} pins no coefficient at the top order;
the little-$o$ refinement does. This tide adds
\code{ForwardExpansionDomain N} as a one-field mixin carrying that
refinement, and proves the exact exponent split it exists to feed:
for $q > 0$ at a critical point,
\[
  \frac{L(qz) - L(0)}{q^2} = T_2(z)
  + \sum_{s=1}^{N} q^s\, V_s(z) + q^N \rho_q(z),
\]
where $V_s = T_{s+2}$ is the degree-$(s{+}2)$ diagonal Taylor term
and the scaled remainder $\rho_q(z) = R(qz)/q^{N+2}$ satisfies the
two estimates the graded exponential stage consumes: pointwise
vanishing $\rho_q(z) \to 0$ as $q \to 0^+$ (from the Peano field
alone), and the uniform window bound
$|\rho_q(z)| \le C\,\|z\|^{N+2}$ wherever $q$-scaling lands in the
Taylor ball (from the quantitative field plus an operator-norm bound
on the top term).
\end{abstract}

\section{Setting}

Stage 3 of the archived implementation order, between the mesoscopic
window (stage 2) and the graded exponential (stage 4). The design
consult was explicit that the mixin should be a single Peano field
over \code{HigherLaplaceDomain (N+2)} rather than a fresh hypothesis
package, and that the split must be an exact algebraic identity in
$q$, not an asymptotic statement: all the analysis lives in
$\rho_q$, and the two lemmas about $\rho_q$ are the entire analytic
content of the stage. The identification of $T_2$ with the quadratic
form of $H$ is deliberately not consumed here; it is bridged
separately, so the split stays a statement about diagonal Taylor
terms only.

\section{The thing formalised}

The mixin \code{ForwardExpansionDomain N L H} extends
\code{HigherLaplaceDomain (N+2)} by the field \code{taylorPeano}:
the degree-$(N{+}2)$ Taylor remainder is $o(\|y\|^{N+2})$ at the
origin. Small general facts: $T_0 = L(0)$
(\code{taylorHomogeneousTerm\_zero}), positive-degree terms vanish
at the origin (via homogeneity at scalar $0$), and the pointwise
operator-norm bound (\code{abs\_taylorHomogeneousTerm\_le})
\[
  |T_k(z)| \le (k!)^{-1}\,\|D^k L(0)\|\,\|z\|^k.
\]
On the mixin:
\code{taylorRem}/\code{scaledRem}, the pointwise vanishing
\code{tendsto\_scaledRem} (compose the Peano little-$o$ along
$q \mapsto qz$, rewrite $\|qz\|^{N+2} = \|z\|^{N+2} q^{N+2}$ on
$q > 0$, absorb the constant, and apply
\code{IsLittleO.tendsto\_div\_nhds\_zero}; the $z = 0$ case is the
remainder vanishing at the origin), the window bound
\code{abs\_scaledRem\_le} with constant
$C = C_{\mathrm{Taylor}} + ((N{+}2)!)^{-1}\|D^{N+2}L(0)\|$
(split the range-$(N{+}3)$ remainder as the domain's range-$(N{+}2)$
remainder minus the top term), and the split
\code{exponent\_split} itself, whose hypothesis is the vanishing of
the degree-one diagonal term --- the critical-point condition in the
language the file already speaks.

\section{Where progress stalled}

Once. The endgame of the split originally routed through a bespoke
helper equating $R/q^{N+2} \cdot q^N \cdot q^2 = R$ by a rewrite
chain, feeding \code{linarith}; the chain broke on association and
\code{linarith} failed on atom mismatch. The fix was to delete the
helper entirely: with $q^2 \ne 0$ and $q^{N+2} \ne 0$ in context,
\code{field\_simp} plus \code{ring} closes the split directly,
because \code{ring} normalizes $q^{N+2} = q^N q^2$ with a variable
$N$. The triple bottom-peel of the sum
(\code{Finset.sum\_range\_succ'} three times) worked first try.

\section{Mathlib lookup versus new mathematics}

\begin{sloppypar}
Lookup: \code{iteratedFDeriv\_zero\_apply}, the multilinear
\code{le\_opNorm}, \code{Finset.sum\_range\_succ'}, and
\code{IsLittleO.comp\_tendsto} with
\code{tendsto\_div\_nhds\_zero}.
New: nothing mathematical --- the split is the note's equation (the
germ of the expansion in the nondegenerate case), stated exactly.
\end{sloppypar}

\section{What the next tide inherits}

Stage 4: the recursive correction coefficients
$P_0 = 1$, $P_j = -\tfrac1j \sum_{s=1}^{j} s\,V_s P_{j-s}$, and the
graded expansion of $\exp(-\sum_s q^s V_s - q^N \rho_q)$ with a
quantitative companion on the mesoscopic window, consuming this
tide's \code{tendsto\_scaledRem} and \code{abs\_scaledRem\_le}
together with stage 2's window geometry. Then the numerator
expansion (stage 5), division (stage 6), and the public existence
and jet-comparison theorems (stage 7).
\end{document}
